% document formatting
\documentclass[10pt]{article}
% \usepackage[UTF8]{inputenc}
\usepackage[UTF8, scheme=plain]{ctex}
\usepackage[left=1in,right=1in,top=1in,bottom=1in]{geometry}
\usepackage[T1]{fontenc}
\usepackage{xcolor}

% math symbols, etc.
\usepackage{amsmath, amsfonts, amssymb, amsthm}

% lists
\usepackage{enumerate}
\usepackage[table,xcdraw]{xcolor}

% images
\usepackage{graphicx} % for images
\usepackage{tikz}

% code blocks
% \usepackage{minted, listings} 

% verbatim greek
\usepackage{alphabeta}

\graphicspath{{./assets/images/Week 2}}

\title{02-741 Week 2 \\ \large{Generative AI for Biomedicine}}
 
\author{Aidan Jan}

\date{\today}

\begin{document}
\maketitle

\section*{Modeling Structure with Graph Neural Networks}
Graphs are everywere!
\begin{itemize}
	\item Molecular graphs
	\item Gene graphs
	\item Brain graphs
	\item Web graphs
	\item Transportation graphs
	\item Social graphs
\end{itemize}
There are a few common tasks we can do on graph data:
\begin{itemize}
	\item Node-level:
	\begin{itemize}
	    \item Link prediction: is there a bond between two atoms?
	    \item Node classification: is this atom C or N?
	    \item Structure prediction: what are 3D coordinates of atoms / residues?
    \end{itemize}
	\item Graph-level:
	\begin{itemize}
	    \item Predicting properties of a graph: Is a molecule toxic to animals?
    \end{itemize}
\end{itemize}

\subsection*{Graph Representation}
\begin{itemize}
	\item Graph: $G = \{V, E\}$
	\item Nodes: $V = \{v_1, v_2, \dots, v_N\}$
	\item Edges: $E = \{e_1, e_2, \dots, e_M\} \subset V \times V$
\end{itemize}

\subsection*{Atom-level Molecule Graph}
\begin{itemize}
	\item Node: heavy (non-hydrogen) atoms
	\item Edge: bond
\end{itemize}
\begin{center} 
	\includegraphics*[width=0.8\textwidth]{W2_1.png} 
\end{center}

\subsubsection*{Node Embedding}
\[Enc(\cdot): V \rightarrow \mathbb{R}^d\]
\begin{center} 
	\includegraphics*[width=0.5\textwidth]{W2_2.png} 
\end{center}
A \textbf{``shallow''} node embedding is just a lookup table.
\begin{center} 
	\includegraphics*[width=0.7\textwidth]{W2_3.png} 
\end{center}

\subsection*{Graph Neural Network}
\begin{itemize}
	\item Input is a graph.
	\item There are many hidden GNN layers.
	\item Output is an embedding matrix for nodes for further downstream tasks: e.g., node property prediction.
\end{itemize}

\subsection*{GNN-based Foundation Models for Biomedicine}
\begin{itemize}
	\item Predicting docking
	\begin{itemize}
	    \item EquiDock
	    \item EquiBind
    \end{itemize}
	\item Property prediction
	\begin{itemize}
	    \item EnviroDataNet
	    \item TorchMD-NET
    \end{itemize}
	\item Modeling protein
	\begin{itemize}
	    \item GearNet
    \end{itemize}
	\item Gene regulator network
	\begin{itemize}
	    \item DeepSEM
    \end{itemize}
	\item Drug design
	\begin{itemize}
	    \item JT-VAE
	    \item GraphAF
    \end{itemize}
\end{itemize}

\section*{Message Passing Neural Network}
Every node receives message from its neighbors, aggregate, and update
\begin{itemize}
	\item $h_i$: node embedding vector.  Aggregates information from its neighbors
	\[h_i^{k + 1} = U^k \left(\text{Aggregate}_{v_j \in N(v_i)} m_{j \rightarrow i}^k\right), \forall v_i \in V\]
	\item $N(v_i)$: Neighbors of the node $V_i$
	\item $m_{j \rightarrow i}^k$: message from node $j$ to $i$ at $k$-th layer.  Aggregate can be average, sum, max/min, etc.
\end{itemize}
\begin{center} 
	\includegraphics*[width=0.5\textwidth]{W2_4.png} 
\end{center}

\subsubsection*{MPNN Message Computation}
\begin{itemize}
	\item Each node will create a message $m_{j \rightarrow i} = M^k (h_i^k, h_j^k, e_{ij})$
	\begin{itemize}
	    \item Dependent on incoming node $(h_j^k)$, target node $(h_i^k)$, edge feature
    \end{itemize}
    \item e.g., Linear projection or FFN $\sigma(W_k \cdot [h_i^k, h_j^k, e_{ij}])$
\end{itemize}
\begin{center} 
	\includegraphics*[width=0.7\textwidth]{W2_5.png} 
\end{center}

\subsubsection*{Message Aggregation / Pooling}
\begin{itemize}
	\item Each node will aggregate messages from its neighbors
	\item Aggregate function can be:
	\begin{itemize}
	    \item Sum, Mean, Max operator
	    \item $m_i^{k + 1} = \sum_{v_j \in N(v_i)}M^k(h_i^k, h_j^k, e_{ij})$
    \end{itemize}
\end{itemize}
\begin{center} 
	\includegraphics*[width=0.7\textwidth]{W2_6.png} 
\end{center}

\subsubsection*{Message Update}
\begin{itemize}
	\item Combine message and current node embedding
	\item Apply FFN
	\begin{itemize}
	    \item $h_i^{k + 1} = U^k (h_i^k, m_i^{k + 1})$
    \end{itemize}
\end{itemize}

\subsubsection*{Generic Framework: Message Passing NN}
\begin{itemize}
	\item MPNN = Message passing + Aggregation + Update
	\begin{itemize}
	    \item propagate nonlinear messages along graph topology
	    \item flexible and generic design choices under this framework
	    \item Graph convolutional network (GCN)
	    \item GraphSAGE
	    \item Graph Attention Network (GAT)
    \end{itemize}
\end{itemize}
\begin{center} 
	\includegraphics*[width=0.9\textwidth]{W2_7.png} 
\end{center}

\subsection*{A Simple Graph Convolutional Network}
\begin{itemize}
	\item Averaging neighbor's message and apply nonlinear transformation
	\begin{itemize}
	    \item Initial embedding: $h_i^0 = x_i$
	    \item Computing next layer: 
        \[h_i^{k + 1} = \sigma \left(W_k \frac{1}{|N(v_i)} \sum_{v_j \in N(v_i)} h_j^k + B_k h_i^k \right)\]
    \end{itemize}
\end{itemize}
Consider the example graph:
\begin{center} 
	\includegraphics*[width=0.5\textwidth]{W2_4.png} 
\end{center}
\begin{itemize}
    \item Example: $h_1^2 = \tanh \left(W_1 \cdot \frac{1}{3} (h_3^1 + h_5^1 + h_8^1) + B_1 h_1^1\right)$
    \item Next layer: $h_1^{(3)} = \tanh \left(W_2 \cdot \frac{1}{3} (h_3^{(2)} + h_5^{(2)} + h_8^{(2)}) + B_2 h_1^{(2)}\right)$
\end{itemize}

\subsection*{Matrix Representations of Graphs}
We typically use an adjacency matrix:
\begin{itemize}
	\item $A[i, j] = 1$ if $v_i$ is adjacent to $v_j$.
	\item $A[i, j] = 0$, otherwise.
	\item Adjacency matrices for undirected graphs are symmetrical over the main diagonal.
\end{itemize}

\subsection*{Matrix Representation of GCN}
\begin{itemize}
	\item Neighbor Aggregation can be performed efficiently using matrix operations
	\[H^k = [h_1^k, \dots, h_{|V|}^k]^T\]
    \item Then,
    \[\sum_{v_j \in N(v_i)} h_j^k = A_{i, :} H^k\]
    \item Let $D$ be diagonal matrix $D_{i, i} = \text{Degree}(v_i) = \sum_j A_{i, j}$
    \begin{itemize}
	    \item The Degree function returns the number of neighbors a given node has.
    \end{itemize}
    \item Then, 
    \[\frac{1}{|N(v_i)|} \sum_{v_j \in N(v_i)} h_j^k = D^{-1} AH^k\]
\end{itemize}

\subsection*{Aggregation of Neighbor's Information in Matrix Form}
Neighbor aggregation can be performed efficiently using matrix operations
\begin{center} 
	\includegraphics*[width=0.9\textwidth]{W2_8.png} 
\end{center}

\subsubsection*{Graph Convolution in Matrix Form}
\begin{align*}
    H^k &= [h_1^k, \dots, h_{|V|}^k]^T \\
    \tilde{A} &= D^{-1} A \\
    H^{k + 1} &= \sigma(\tilde{A} H^k \cdot W_k^T + H^k B_k^T) \\
    h_i^{k + 1} &= \sigma \left(W_k \frac{1}{|N(v_i)|} \sum_{v_j \in N(v_i)}h_j^k + B_k h_i^k\right)
\end{align*}

\subsection*{(Better) Graph Convolutional Network}
To make $\tilde{A}$ symmetric, we keep the definition of $H^k$:
\[H^k = [h_1^k, \dots, h_{|V|}^k]^T\]
But, 
\[\tilde{A} = D^{-\frac{1}{2}} A D^{-\frac{1}{2}}\]
We also keep the same definition of $H^{k + 1}$:
\[H^{k + 1} = \sigma(\tilde{A} H^k \cdot W_k^T + H^k B_k^T)\]

\subsection*{Prediction Layer}
To predict, we apply softmax functions
\begin{itemize}
	\item For node classification, (e.g. predicting atom type)
	\[o_i = \text{Softmax}(h_i^{(m)})\]
    \item For graph classification, (e.g., predicting toxicity)
    \[o = \text{Softmax}(\frac{1}{N} \sum_i h_i^{(m)})\]
\end{itemize}

\subsection*{Property of GCN: Equivariant}
The embeddings computed from graph convolution layers is \textbf{invariant} to node order permutation.

\subsection*{Training GCN}
\begin{itemize}
	\item Parameters: weight matrix for each layer
	\item Supervised training: e.g., Node classification
	\begin{itemize}
	    \item Cross entropy loss
	    \[L = -\sum y_i \log f(h_i^K)\]
        \[f_i = \text{Softmax}(h_i^{(K)})\]
        \item $y_i$ is node label (in one-hot vector or smoothed version)
    \end{itemize}
    \item Unsupervised training:
    \begin{itemize}
	    \item Linked nodes have similar embedding
	    \[L = \sum_{i, j} CE(y_{i, j}, Sim(h_i^K, h_j^K))\]
        \item $y_{i, j} = 1$ if there is an edge from $v_i$ to $v_j$
        \item Similarity can be defined in many ways: e.g., inner product $h_i \cdot h_j$
    \end{itemize}
\end{itemize}

\subsection*{GraphSAGE}
Also a special case of MPNN
\[h_i^{k + 1} = \sigma(W_k \cdot \text{CONCAT}(h_i^k, \text{AGG}(\{h_j^k, \forall v_j \in N(v_i)\})))\]
AGG is the aggregate function and can be designed in multiple ways, like pooling (sum, avg, max).

\subsection*{Graph Attention Network (GAT)}
Not strictly MPNN, but a generalized one:
\[h_i^{k + 1} = \sigma \left(\sum_{v_j \in N(v_i)} \alpha_{ij} W_k h_{v_j}^k\right)\]
Attention weights:
\[\alpha_{ij} = \text{Attention}(W_k h_i, W_k h_j) = \frac{\exp(W_h h_i)^T W_k h_j}{\sum_{j'} \exp(W_k h_i)^T W_k h_{j'}}\]
Multi-head attention can also be used on GAT.

\subsection*{Relation between GNN and CNN}
A CNN (convolutional neural network) can be viewed as a special GNN on a grid graph.  Essentially, one node has 8 nodes surrounding it, four adjacent, four diagonal.

\subsection*{GNN vs. Transformer}
A transformer is a special GNN on a fully-connected graph.

\section*{Equivariant Graph Neural Network}
\subsection*{Structure and Geometry of Molecules}
\begin{itemize}
	\item Matrix of 3D coordinates, in addition to node features
	\item Equivariance: translation, rotation invariant
\end{itemize}

\subsection*{Equivariance in Embedding Transformation}
\begin{itemize}
	\item $f$ is Equivariant (SE3 invariant):
	\[f(R(x) + z) = f(x)\]
    \item No matter how we rotate and move the molecule, the NN will compute the same embeddings.
    \item $X$ are coordinates
    \item $H$ are features
    \item $g$ is the rotation / translation.
    \[F(H, g(X)) = F(H, X)\]
\end{itemize}
\begin{center} 
	\includegraphics*[width=0.7\textwidth]{W2_9.png} 
\end{center}
A \textbf{Equivariant Graph Neural Network (EGNN)} is an equivariant MPNN, optionally updating coordinates.

\begin{align*}
    d_{ij} &= |x_i - x_j|_2 \\
    m_{j \rightarrow i} &= M^k (h_i^k, h_j^k, d_{ij}) \\
    m_i^{k + 1} &= \sum_{v_j \in N(v_i)} m_{j \rightarrow i} \\
    h_i^{k + 1} &= U^k (h_i^k, m_i^{k + 1})
\end{align*}

\section*{Auto-Encoder}
Learning representations so that data can be reconstructed!
[FILL 4]
Need to avoid learning an identity mapping
\begin{itemize}
	\item Simple instance:
	\begin{itemize}
	    \item encoder is linear projection
	    \item decoder is linear projection
	    \item weight typing
    \end{itemize}
	\item Other constraints:
	\begin{itemize}
	    \item Robust construction to noise, (Denoising Auto-Encoder)
	    \item Hidden representations conform to a distribution (Variational Auto-Encoder)
    \end{itemize}
\end{itemize}
[Fill 5]
The box with $h$ is a probability function (to provide the variability for the Variational Auto-Encoder)

\subsection*{Training Auto-Encoder}
\begin{itemize}
	\item Real-valued data (e.g., images)
	\begin{itemize}
	    \item MSE loss
	    \item $\text{MSE} = ||f(x) - x||$
    \end{itemize}
    \item Categorical data (e.g., sentences)
    \begin{itemize}
	    \item Cross-Entropy
	    \item $\text{CE} = -\sum x_i \log f(x)_i$
    \end{itemize}
\end{itemize}

\subsection*{Denoising Auto-Encoder}
\begin{itemize}
	\item Encoder takes a noised data to produce the hidden representation
	\item Decoder needs to recover original data from hidden representation
	\item Make the auto-encoder robust
\end{itemize}
[Fill 9]

\subsection*{Ovoercomplete Representation}
\begin{itemize}
	\item If hidden representation is larger than input data dimension
	\item Need to enforce additional constraints (e.g., sparsity)
	\begin{itemize}
	    \item [FILL loss]  
    \end{itemize}
\end{itemize}

\subsection*{Variational Auto-Encoder (VAE)}
\begin{itemize}
	\item Input, intermediate, and output values are no longer deterministic vectors, but are probabilistic variables
	\item Using probabilistic graphical model (PRG) to represent dependencies
\end{itemize}

\section*{Probabilistic Graphical Model}
\begin{itemize}
	\item Models describe how the world works
	\item Models are always simplifications
	\begin{itemize}
	    \item May not account for every variable
	    \item May not account for all interactions between variables
    \end{itemize}
    \item We=hat do we do with probabilistic models?
    \begin{itemize}
	    \item We (or our agents) need to reason about unknown variables, given evidence
	    \item Example: explanation (diagnostic reasoning)
	    \item Example: prediction (causal reasoning)
	    \item Example: value of information
    \end{itemize}
\end{itemize}

\subsection*{Directed PGM - Bayesian Networks}
\begin{itemize}
	\item One type of PGM with directed graph
	\item Nodes: random variables (with domains)
	\begin{itemize}
	    \item Can be assigned (observed) or unassugned (unobserved)
    \end{itemize}
    \item Arcs: interactions
    \begin{itemize}
	    \item Indicate ``direct influence'' between variables
	    \item Formally: encode conditional independence (more later)
    \end{itemize}
    \item For now: imagine that arrows mean direct causation (in general, they don't!)
\end{itemize}

\subsection*{Bayesian Network Semantics}
\begin{itemize}
	\item A set of nodes, one per variable $X$
	\item A directed, acyclic graph (DAG)
	\item A conditional distribution for each node
	\begin{itemize}
	    \item A collection of distributions ovver $X$, one for each combination of parents' values
	    \[P(X | a_1, \dots, a_n)\]
        \item CPT: conditional probability table
        \item Description of a noisy ``causal'' table
    \end{itemize}
    \item Directed graphical models
\end{itemize}

\subsubsection*{Example: Rain and Traffic}
\begin{itemize}
	\item Variables: 
	\begin{itemize}
	    \item $R$: if it rains (1) otherwise (0)
	    \item $T$: if there is traffic (1) otherwise (0)
    \end{itemize}
    \item Model 1: independence.
    \item Model 2: rain causes traffic
    [FILL tables + image]
\end{itemize}

\subsubsection*{Example: Traffic II}
\begin{itemize}
	\item Variables:
	\begin{itemize}
	    \item T: Traffic
	    \item R: It rains
	    \item U: Umbrella
    \end{itemize}
    \item Q: What is the probability of rain if we observe traffic but no umbrella?
    \begin{itemize}
	    \item $P(R = 1 | T = 1, U = 0)$
    \end{itemize}
\end{itemize}

\section*{Variational Auto-Encoders}
[FILL (before 20)]
\subsection*{Deep Latent Model}
\begin{itemize}
	\item $z$ follows a prior distribution, e.g., Gaussian(0, I)
	\item $p(x | z)$ is defined by a deep neural network $f(z; \theta)$
	\begin{itemize}
	    \item We use a deep neural network to simulate the probability functions.
    \end{itemize}
	\item To learn $\theta$, use $E_{(Z|X)} [\log p(X, Z;\theta)]$
\end{itemize}

\subsection*{Probabilistic Graphical Model for VAE}
\begin{itemize}
	\item Assuming data $X$ is generated from a latent variable $Z$
	\item Generation process:
	\begin{itemize}
	    \item draw $Z \sim N(\mu, \Sigma)$
	    \item draw $X | Z \sim p(f(Z))$, defined by a neural network $f$ with parameter $\sigma$
    \end{itemize}
    \begin{center} 
        [FILL 21]
    \end{center}
\end{itemize}

\subsection*{Derivation of VAE}
Objective: maximize the data log-likelihood
\begin{align*}
    \max \ell(\theta) &= \sum_n \log p(x_n; \theta) \\
    &= \sum_n \log \int p(x_n | z_n; \theta) p(z_n; \theta) \text{ d} z_n
\end{align*}
\begin{itemize}
	\item Hard to optimize over $\theta$, if $f(Z)$ is very complex, due to the integration in the log.  (i.e., $p(x_n | z_n; \theta)$ is very complex)
	\item Examples of complex $f(Z)$ include CNNs, RNNs, GNNs, Transformers
\end{itemize}

\subsubsection*{Variational Posterior Distribution}
\begin{itemize}
	\item See equations above
	\item $\log p(x; \theta)$ is intractable.
	\item For any distribution $q(z | x, \phi)$:
	\[\log p(x; \theta) \geq E_{q(z|x; \phi)} \left[\log \frac{p(x, z; \theta)}{q(z|x; \phi)}\right] = \text{ELBO}\]
    \item Derivation via Jensen's inequality.
    \begin{itemize}
	    \item If we want to find $f(E(x))$, where $E(x)$ is an expectation (e.g., average value), and $f$ is a \textit{concave} function, then $E[f(x)] < f(E(x))$.
	    \item It turns out in the ELBO equation, we are essentially maximizing a lower bound for our equation.
    \end{itemize}
    \item Maximizing the ELBO instead of maximizing $\log p(x; \theta)$
\end{itemize}

\subsubsection*{Understanding ELBO}
\begin{align*}
    \log p(X; \theta) &\geq E_q\left[\log \frac{p(X, Z; \theta)}{q(Z | X; \phi)}\right] \\
    &= \sum_n E_q \left[\log \frac{p(x_n | z_n; \theta) p_0(z_n)}{q(z_n | x_n; \phi)}\right] \\
    &= E_q [\log p(x_n | z_n; \theta)] - KL(q(z_n|x_n; \phi) || p_0(z_n))
\end{align*}
\begin{itemize}
	\item The first term (before subtraction) is reconstruction loss
	\item The second term is regularization.  (KL = Kullback-Leibler divergence)
\end{itemize}

\subsection*{Variational Auto-Encoder with Neural Network}
Let $q(z | x; \phi)$ and $p(x | z; \theta)$ share the same parameter $\theta$
[FILL 25]

\subsection*{Training VAE}
\begin{itemize}
	\item Gradient Descent
	\item [FILL]
\end{itemize}

[FILL]

\subsection*{VAE for Sequence Data}



\end{document}